%% Template for a new course sheet.
%% Copy this to <course number>.tex, e.g. 4551.tex, and fill it in.

\documentclass{osucourse}
\course{4551}          % <-- your course number, matching the file name

\begin{document}

%% The catalog description, prerequisites and exclusions are filled in from
%% the university's data.  Leave these two lines alone.
\CatalogDescription
\PrerequisitesAndExclusions

%% Delete any section below that your course does not need.

\begin{Purpose}
What this course is for, and how it fits into the major.
\end{Purpose}

\begin{LearningOutcomes}
By the end of this course, students should be able to:
\begin{itemize}
\item do the first thing;
\item do the second thing.
\end{itemize}
\end{LearningOutcomes}

%% Only the 1000-level algebra and precalculus sheets show this; it is one
%% shared figure, in sequencing-chart.tex.
%\SequencingChart

\begin{Text}
\emph{Title of the Book}, 5\textsuperscript{th} edition, by Author,
published by Publisher, ISBN: 9780000000000
\end{Text}

\begin{Technology}
Any calculator or software the course requires.
\end{Technology}

\begin{TopicsList}
\item First topic.
\item Second topic.
\end{TopicsList}

%% If your topics list is an outline rather than a flat numbered list --
%% chapter headings, numbering that restarts -- use TopicsOutline instead:
%%
%%   \begin{TopicsOutline}
%%   \textbf{1. Preliminaries}
%%   1.1 Sets and equivalence relations
%%   \end{TopicsOutline}

\begin{SuggestedSchedule}
\week{1}{What happens in week one.}
\week{2}{What happens in week two. \emph{Midterm I.}}
\finalsweek{Final examination.}
\end{SuggestedSchedule}

\end{document}
